% answers.tex —— 1990 年·广东卷·理科 · 参考答案与详细解析片段
% 由 tools/generate_projects.py 依来源总账自 1990/1990_数学真题_广东理_试题.md + 1990/1990_数学真题_广东理_解析正文.md + 1990/1990_数学真题_广东理_答案要点.md 生成，勿手改（改源 MD 后重跑）


\begin{gktip}
第一卷已完成独立盲核，第二卷来源解答与主会话整理内容仍分开标识。第 2 题按通常定义 A、C 同时成立，与单选形式及来源表仅印 A 冲突，未裁定唯一答案；第 12 题原卷图位缺失，虽来源表记 B，仍不能核验数学结论。
\end{gktip}

\begin{answer}
\gksolhead{1}{}

\textbf{题目}：

如果直线 $ax+2y+1=0$ 与直线 $x+y-2=0$ 互相垂直，那么 $a$ 的值等于（　）。

\begin{gkopts}(4)
\task A．$1$
\task B．$-\dfrac{1}{3}$
\task C．$-\dfrac{2}{3}$
\task D．$-2$
\end{gkopts}

\Needspace{4\baselineskip}
\textbf{解析}

两直线斜率分别为 $-\dfrac a2$ 和 $-1$。垂直时斜率乘积为 $-1$，故 $\dfrac a2=-1$，得 $a=-2$。

\gkanswer{答案：D}

\end{answer}
\begin{answer}
\gksolhead{2}{}

\textbf{题目}：

如果直线 $l$ 是平面 $\alpha$ 的斜线，那么在平面 $\alpha$ 内（　）。

\begin{gkopts}(1)
\task A．不存在与 $l$ 平行的直线；B．不存在与 $l$ 垂直的直线；
\task C．与 $l$ 垂直的直线只有一条；D．与 $l$ 平行的直线有无穷多条。
\end{gkopts}

\Needspace{4\baselineskip}
\textbf{解析}

斜线 $l$ 与平面 $\alpha$ 交于一点 $P$，其在 $\alpha$ 内的正投影方向非零。平面内过 $P$ 且垂直于该投影方向的直线只有一条；按通常“垂直”定义（两直线相交且成直角），这条直线也垂直于 $l$，所以 C 也成立。另一方面，斜线方向含有非零的平面法向分量，故平面内不存在与 $l$ 平行的直线，A 同样成立。于是按通常定义A、C均成立，与单选题形式及来源答案表仅印 A 不相容；在取得勘误或可靠原题说明前，保留题面和印载答案，不把本题判作唯一答案。

\gkanswer{答案：A}

\end{answer}
\begin{answer}
\gksolhead{3}{}

\textbf{题目}：

某乒乓球队有 $9$ 名队员，其中 $2$ 名是种子选手，现要挑选 $5$ 名队员参加比赛，种子选手都必须在内，那么不同的选法共有（　）。

\begin{gkopts}(4)
\task A．$126$ 种
\task B．$84$ 种
\task C．$35$ 种
\task D．$21$ 种
\end{gkopts}

\Needspace{4\baselineskip}
\textbf{解析}

两名种子选手已固定入选，只需从其余 $7$ 人中选 $3$ 人，方法数为 $\binom73=35$。

\gkanswer{答案：C}

\end{answer}
\begin{answer}
\gksolhead{4}{}

\textbf{题目}：

把 $6$ 个不同的元素排成前后两排，每排 $3$ 个元素，那么不同的排法共有（　）。

\begin{gkopts}(4)
\task A．$36$ 种
\task B．$120$ 种
\task C．$720$ 种
\task D．$1440$ 种
\end{gkopts}

\Needspace{4\baselineskip}
\textbf{解析}

前后两排共占 $6$ 个有区别的位置，$6$ 个不同元素的排法数为 $6!=720$。

\gkanswer{答案：C}

\end{answer}
\begin{answer}
\gksolhead{5}{}

\textbf{题目}：

一个圆台的母线长是上、下底面半径的等差中项，且侧面积为 $8\pi\,\mathrm{cm}^2$，那么母线长为（　）。

\begin{gkopts}(2)
\task A．$4\,\mathrm{cm}$
\task B．$2\sqrt2\,\mathrm{cm}$
\task C．$2\,\mathrm{cm}$
\task D．$\sqrt2\,\mathrm{cm}$
\end{gkopts}

\Needspace{4\baselineskip}
\textbf{解析}

设上、下底面半径为 $R,r$，母线长为 $l$。题设给出 $R+r=2l$；圆台侧面积为 $\pi(R+r)l=8\pi$，所以 $2\pi l^2=8\pi$，得 $l=2\,\mathrm{cm}$。

\gkanswer{答案：C}

\end{answer}
\begin{answer}
\gksolhead{6}{}

\textbf{题目}：

如果复数 $z=3+ai$ 满足条件 $|z-2|<2$，那么实数 $a$ 的取值范围是（　）。

\begin{gkopts}(2)
\task A．$(-2\sqrt2,2\sqrt2)$
\task B．$(-2,2)$
\task C．$(-1,1)$
\task D．$(-\sqrt3,\sqrt3)$
\end{gkopts}

\Needspace{4\baselineskip}
\textbf{解析}

由 $z=3+ai$ 得 $|z-2|=|1+ai|=\sqrt{1+a^2}<2$，故 $a^2<3$，即 $-\sqrt3<a<\sqrt3$。

\gkanswer{答案：D}

\end{answer}
\begin{answer}
\gksolhead{7}{}

\textbf{题目}：

一动点在圆 $x^2+y^2=1$ 上移动时，它与定点 $(3,0)$ 连线的中点的轨迹方程是（　）。

\begin{gkopts}(2)
\task A．$(x+3)^2+y^2=4$
\task B．$(x-3)^2+y^2=1$；
\task C．$(2x-3)^2+4y^2=1$
\task D．$\left(x+\dfrac32\right)^2+y^2=\dfrac12$。
\end{gkopts}

\Needspace{4\baselineskip}
\textbf{解析}

设圆上动点为 $(u,v)$，中点为 $(x,y)$，则 $u=2x-3$、$v=2y$。代入 $u^2+v^2=1$，得 $(2x-3)^2+4y^2=1$。

\gkanswer{答案：C}

\end{answer}
\begin{answer}
\gksolhead{8}{}

\textbf{题目}：

\Needspace{5\baselineskip}
函数
\nopagebreak[4]
\[
y=\frac{\sin x}{|\sin x|}+\frac{|\cos x|}{\cos x}+\frac{\operatorname{tg}x}{|\operatorname{tg}x|}+\frac{|\operatorname{ctg}x|}{\operatorname{ctg}x}
\]
的值域是（　）。

\begin{gkopts}(4)
\task A．$\{-2,4\}$
\task B．$\{-2,0,4\}$
\task C．$\{-2,0,2,4\}$
\task D．$\{-4,-2,0,4\}$
\end{gkopts}

\Needspace{4\baselineskip}
\textbf{解析}

令 $s=\operatorname{sgn}(\sin x)$、$c=\operatorname{sgn}(\cos x)$。题中四项之和为 $s+c+2sc$；当 $s,c$ 取 $\pm1$ 时，所得值为 $4,-2,0$，故值域为 $\{-2,0,4\}$。

\gkanswer{答案：B}

\end{answer}
\begin{answer}
\gksolhead{9}{}

\textbf{题目}：

方程 $\sin 2x=\sin x$ 在区间 $(0,2\pi)$ 内的解的个数是（　）。

\begin{gkopts}(4)
\task A．$4$
\task B．$3$
\task C．$2$
\task D．$1$
\end{gkopts}

\Needspace{4\baselineskip}
\textbf{解析}

方程化为 $\sin x(2\cos x-1)=0$。在 $(0,2\pi)$ 内，$\sin x=0$ 给 $x=\pi$；$\cos x=\dfrac12$ 给 $x=\dfrac\pi3,\dfrac{5\pi}3$，共 $3$ 个解。

\gkanswer{答案：B}

\end{answer}
\begin{answer}
\gksolhead{10}{}

\textbf{题目}：

函数 $y=(x+4)^2$ 在某区间上是减函数，这区间可以是（　）。

\begin{gkopts}(4)
\task A．$(-\infty,-4]$
\task B．$[-4,+\infty)$
\task C．$[4,+\infty)$
\task D．$(-\infty,4]$
\end{gkopts}

\Needspace{4\baselineskip}
\textbf{解析}

抛物线 $y=(x+4)^2$ 的对称轴为 $x=-4$，在 $x\le-4$ 时递减，故可取区间 $(-\infty,-4]$。

\gkanswer{答案：A}

\end{answer}
\begin{answer}
\gksolhead{11}{}

\textbf{题目}：

已知函数：① $y=\operatorname{arctg}x$，② $y=\dfrac\pi2-\operatorname{arcctg}x$。那么（　）。

\begin{gkopts}(1)
\task A．①与②都是奇函数；B．①与②都是偶函数；
\task C．①是奇函数，②是偶函数；D．①与②都不是奇函数，也都不是偶函数。
\end{gkopts}

\Needspace{4\baselineskip}
\textbf{解析}

反正切函数为奇函数。按反余切函数通常取值范围 $(0,\pi)$，有 $\operatorname{arcctg}(-x)=\pi-\operatorname{arcctg}x$，于是 $\dfrac\pi2-\operatorname{arcctg}x$ 也为奇函数。

\gkanswer{答案：A}

\end{answer}
\begin{answer}
\gksolhead{12}{}

\textbf{题目}：

已知下图是函数 $y=2\sin(\omega x+\varphi)$（$|\varphi|<\dfrac\pi2$）的图象，那么（　）。

\begin{gktip}
\textbf{缺源提示：}固定来源 PDF 第 2 页该题图位为空白。以下只照录题干和选项；不依据答案表重建图形，题目保持待核。
\end{gktip}

\begin{gkopts}(1)
\task A．$\omega=2,\ \varphi=-\dfrac\pi6$
\task B．$\omega=2,\ \varphi=\dfrac\pi6$；
\task C．$\omega=\dfrac{10}{11},\ \varphi=\dfrac\pi6$
\task D．$\omega=\dfrac{10}{11},\ \varphi=-\dfrac\pi6$。
\end{gkopts}

\Needspace{4\baselineskip}
\textbf{解析}

固定来源 PDF 第 2 页的题图缺失，无法依据题面确定 $\omega,\varphi$。来源答案表记录 B，但该记录不能替代缺失的图形条件。

\gkanswer{答案：B（来源答案表记录；题图缺失，待核）}

\end{answer}
\begin{answer}
\gksolhead{13}{}

\textbf{题目}：

如果直线 $y=ax+2$ 与直线 $y=3x-b$ 关于直线 $y=x$ 对称，那么（　）。

\begin{gkopts}(2)
\task A．$a=\dfrac13,\ b=-6$
\task B．$a=\dfrac13,\ b=6$；
\task C．$a=3,\ b=-2$
\task D．$a=3,\ b=6$。
\end{gkopts}

\Needspace{4\baselineskip}
\textbf{解析}

把 $y=ax+2$ 关于 $y=x$ 对称，交换 $x,y$ 得 $x=ay+2$，即 $y=\dfrac1a x-\dfrac2a$。与 $y=3x-b$ 比较，得 $a=\dfrac13$、$b=6$。

\gkanswer{答案：B}

\end{answer}
\begin{answer}
\gksolhead{14}{}

\textbf{题目}：

如果抛物线 $y^2=a(x+1)$ 的准线方程是 $x=-3$，那么它的焦点坐标是（　）。

\begin{gkopts}(4)
\task A．$(3,0)$
\task B．$(2,0)$
\task C．$(1,0)$
\task D．$(-1,0)$
\end{gkopts}

\Needspace{4\baselineskip}
\textbf{解析}

抛物线顶点为 $(-1,0)$。写成标准式 $y^2=4p(x+1)$，其准线为 $x=-1-p=-3$，所以 $p=2$，焦点为 $(1,0)$。

\gkanswer{答案：C}

\end{answer}
\begin{answer}
\gksolhead{15}{}

\textbf{题目}：

如果 $\log_a2>\log_b2>0$，那么（　）。

\begin{gkopts}(4)
\task A．$1<a<b$
\task B．$1<b<a$
\task C．$0<a<b<1$
\task D．$0<b<a<1$
\end{gkopts}

\Needspace{4\baselineskip}
\textbf{解析}

由 $\log_b 2>0$ 得 $b>1$；又因 $\log_a 2>\log_b 2>0$，故 $a>1$。在 $t>1$ 上，$\log_t 2=\dfrac{\ln2}{\ln t}$ 随 $t$ 增大而减小，因此 $\log_a 2>\log_b 2$ 等价于 $a<b$，即 $1<a<b$，对应答案 A。

\gkanswer{答案：A}

\end{answer}
\begin{answer}
\gksolhead{16}{}

\textbf{题目}：

复数 $1+i\operatorname{ctg}\theta$（$\pi<\theta<2\pi$）的三角形式是（　）。

\begin{gkopts}(1)
\task A．$\dfrac1{\sin\theta}(\sin\theta+i\cos\theta)$；
\task B．$\dfrac1{\sin\theta}\left[\cos\left(\dfrac\pi2-\theta\right)+i\sin\left(\dfrac\pi2-\theta\right)\right]$；
\task C．$-\dfrac1{\sin\theta}\left[\cos\left(\dfrac\pi2+\theta\right)+i\sin\left(\dfrac\pi2+\theta\right)\right]$；
\task D．$-\dfrac1{\sin\theta}\left[\cos\left(\dfrac{3\pi}2-\theta\right)+i\sin\left(\dfrac{3\pi}2-\theta\right)\right]$。
\end{gkopts}

\Needspace{4\baselineskip}
\textbf{解析}

有 $1+i\operatorname{ctg}\theta=\dfrac{\sin\theta+i\cos\theta}{\sin\theta}$。因 $\pi<\theta<2\pi$，$\sin\theta<0$；且 $\sin\theta=\cos(\frac\pi2-\theta)$、$\cos\theta=\sin(\frac\pi2-\theta)$。将负振幅改写成正模长后，相位加 $\pi$，得 $\frac{3\pi}2-\theta$，对应 D。

\gkanswer{答案：D}

\end{answer}
\begin{answer}
\gksolhead{17}{}

\textbf{题目}：

如果 $x,y$ 是实数，那么“$xy>0$”是“$|x+y|=|x|+|y|$”的（　）。

\begin{gkopts}(2)
\task A．充分但不必要条件；B．必要但不充分条件；
\task C．必要而且充分的条件；D．既不充分也不必要的条件。
\end{gkopts}

\Needspace{4\baselineskip}
\textbf{解析}

等式 $|x+y|=|x|+|y|$ 等价于 $xy\ge0$。所以 $xy>0$ 是充分条件；但 $x=0$ 或 $y=0$ 时等式仍成立，故它不是必要条件。

\gkanswer{答案：A}

\end{answer}
\begin{answer}
\gksolhead{18}{}

\textbf{题目}：

极坐标方程 $4\rho\sin^2\dfrac\theta2=5$ 表示的曲线是（　）。

\begin{gkopts}(4)
\task A．圆
\task B．椭圆
\task C．抛物线
\task D．双曲线的一支
\end{gkopts}

\Needspace{4\baselineskip}
\textbf{解析}

利用 $\sin^2\dfrac\theta2=\dfrac{1-\cos\theta}{2}$，方程化为 $\rho-x=\dfrac52$，即 $\rho=x+\dfrac52$。平方并用 $\rho^2=x^2+y^2$，得 $y^2=5x+\dfrac{25}{4}$，为抛物线。

\gkanswer{答案：C}

\end{answer}
\begin{answer}
\gksolhead{19}{}

\textbf{题目}：

已知等比数列的公比为 $2$，且前 $4$ 项之和等于 $1$，那么前 $8$ 项之和等于（　）。

\begin{gkopts}(4)
\task A．$15$
\task B．$17$
\task C．$19$
\task D．$21$
\end{gkopts}

\Needspace{4\baselineskip}
\textbf{解析}

公比为 $2$ 时，$S_4=a_1(2^4-1)$、$S_8=a_1(2^8-1)$，故 $\dfrac{S_8}{S_4}=\dfrac{255}{15}=17$。

\gkanswer{答案：B}

\end{answer}
\begin{answer}
\gksolhead{20}{}

\textbf{题目}：

函数 $y=\operatorname{tg}\dfrac x2-\dfrac1{\sin x}$ 的最小正周期是（　）。

\begin{gkopts}(4)
\task A．$\dfrac\pi2$
\task B．$\pi$
\task C．$\dfrac{3\pi}2$
\task D．$2\pi$
\end{gkopts}

\Needspace{4\baselineskip}
\textbf{解析}

在原函数定义域内，$\operatorname{tg}\dfrac x2-\dfrac1{\sin x}=\dfrac{1-\cos x-1}{\sin x}=-\operatorname{ctg}x$。该函数的最小正周期为 $\pi$。

\gkanswer{答案：B}

\end{answer}
\begin{answer}
\gksolhead{21}{}

\textbf{题目}：

$\operatorname{tg}70^\circ+\operatorname{tg}50^\circ-\sqrt3\,\operatorname{tg}50^\circ\operatorname{tg}70^\circ$ 的值等于（　）。

\begin{gkopts}(2)
\task A．$\sqrt3$
\task B．$\dfrac{\sqrt3}3$
\task C．$-\dfrac{\sqrt3}3$
\task D．$-\sqrt3$
\end{gkopts}

\Needspace{4\baselineskip}
\textbf{解析}

由正切和角公式，$\operatorname{tg}70^\circ+\operatorname{tg}50^\circ=-\sqrt3\left(1-\operatorname{tg}70^\circ\operatorname{tg}50^\circ\right)$，故题设表达式等于 $-\sqrt3$。

\gkanswer{答案：D}

\end{answer}
\begin{answer}
\gksolhead{22}{}

\textbf{题目}：

如果实数 $x,y$ 满足等于 $(x-2)^2+y^2=3$，那么 $\dfrac yx$ 的最大值是（　）。

\begin{gkopts}(2)
\task A．$\dfrac12$
\task B．$\sqrt3$
\task C．$\dfrac{\sqrt3}2$
\task D．$\dfrac{\sqrt3}3$
\end{gkopts}

\Needspace{4\baselineskip}
\textbf{解析}

圆心为 $(2,0)$、半径为 $\sqrt3$。从原点引斜率为 $m>0$ 的直线 $y=mx$；当其与圆相切时，圆心到直线距离满足 $\dfrac{2m}{\sqrt{m^2+1}}=\sqrt3$，解得 $m=\sqrt3$。切线斜率给出 $y/x$ 的最大值。

\gkanswer{答案：B}

\end{answer}
\begin{answer}
\gksolhead{23}{}

\textbf{题目}：

以一个正方体的顶点为顶点的四面体共有（　）。

\begin{gkopts}(4)
\task A．$70$ 个
\task B．$64$ 个
\task C．$58$ 个
\task D．$52$ 个
\end{gkopts}

\Needspace{4\baselineskip}
\textbf{解析}

从 $8$ 个顶点中任取 $4$ 个共有 $\binom84=70$ 组。共面四点组对应正方体的 $6$ 个面及 $6$ 个对角截面，共 $12$ 组不构成四面体，故共有 $70-12=58$ 个。

\gkanswer{答案：C}

\end{answer}
\begin{answer}
\gksolhead{24}{}

\textbf{题目}：

设 $z$ 是复数，$z+2$ 的辐角为 $\dfrac\pi3$，$z-2$ 的辐角为 $\dfrac{5\pi}6$，那么 $z$ 等于（　）。

\begin{gkopts}(2)
\task A．$-\dfrac12+\dfrac{\sqrt3}2i$
\task B．$-\dfrac{\sqrt3}2+\dfrac12i$；
\task C．$-\sqrt3+i$
\task D．$-1+\sqrt3i$
\end{gkopts}

\Needspace{4\baselineskip}
\textbf{解析}

设 $z=x+iy$。由两条射线的方向，$y=\sqrt3(x+2)$ 且 $y=-\dfrac{x-2}{\sqrt3}$。联立得 $x=-1,y=\sqrt3$，所以 $z=-1+\sqrt3i$。

\gkanswer{答案：D}

\end{answer}
\begin{answer}
\gksolhead{25}{}

\textbf{题目}：

设函数 $y=\operatorname{arctg}x$ 的图象沿 $x$ 轴正方向平移 $2$ 个单位，所得图象为 $C$。又设图象 $C_1$ 与 $C$ 关于原点对称，那么 $C_1$ 所对应的函数是（　）。

\begin{gkopts}(2)
\task A．$y=-\operatorname{arctg}(x-2)$
\task B．$y=\operatorname{arctg}(x-2)$；
\task C．$y=-\operatorname{arctg}(x+2)$
\task D．$y=\operatorname{arctg}(x+2)$。
\end{gkopts}

\Needspace{4\baselineskip}
\textbf{解析}

向右平移 $2$ 个单位后，$C$ 的函数为 $y=\operatorname{arctg}(x-2)$。关于原点对称后，函数为 $y=-\operatorname{arctg}(-x-2)=\operatorname{arctg}(x+2)$。

\gkanswer{答案：D}

\end{answer}
\begin{answer}
\gksolheadwithlabel{26}{一}{}

\textbf{题目}：

（1） 双曲线 $\dfrac{x^2}{16}-\dfrac{y^2}{9}=1$ 的准线方程是\gkblank{}。

（2） $(x-1)-(x-1)^2+(x-1)^3-(x-1)^4$ 的展开式中，$x^2$ 的系数等于\gkblank{}。

\Needspace{5\baselineskip}
（3） 已知数列 $\{a_n\}$ 是公差不为零的等差数列，如果 $S_n$ 是 $\{a_n\}$ 的前 $n$ 项的和，那么
\nopagebreak[4]
\[
\lim_{n\to\infty}\frac{na_n}{S_n}
\]
的值等于\gkblank{}。

（4） 函数 $y=\sin x+\cos x$ 的最大值是\gkblank{}。

（5） 如图，三棱柱 $ABC-A_1B_1C_1$ 中，若 $E,F$ 分别为 $AB,AC$ 的中点，平面 $EB_1C_1F$ 将三棱柱分成体积为 $V_1,V_2$ 的两部分，那么 $V_1:V_2$ 等于\gkblank{}。
\nopagebreak[4]
\begin{center}\begin{tikzpicture}[x=0.78cm,y=0.66cm,line cap=round,line join=round]
\coordinate (A) at (0,0);
\coordinate (B) at (4,0);
\coordinate (C) at (1.4,1.9);
\coordinate (Aone) at (-0.35,2.0);
\coordinate (Bone) at (3.65,2.0);
\coordinate (Cone) at (1.05,3.9);
\coordinate (E) at (2,0);
\coordinate (F) at (0.7,0.95);

% 原卷三棱柱及平面 EB_1C_1F 的拓扑关系；投影不表示长度比例。
\fill[gray!14] (E)--(Bone)--(Cone)--(F)--cycle;
\draw[thick] (A)--(B)--(C)--cycle;
\draw[thick] (Aone)--(Bone)--(Cone)--(Aone);
\draw[thick] (A)--(Aone) (B)--(Bone) (C)--(Cone);
\draw[dashed] (A)--(C);
\draw[very thick] (E)--(Bone) (Cone)--(F) (F)--(E) (Bone)--(Cone);

\node[below left] at (A) {$A$};
\node[below right] at (B) {$B$};
\node[above right] at (C) {$C$};
\node[above left] at (Aone) {$A_1$};
\node[above right] at (Bone) {$B_1$};
\node[above] at (Cone) {$C_1$};
\node[below] at (E) {$E$};
\node[left] at (F) {$F$};
\node at (0.3,1.55) {$V_1$};
\node at (3.4,1.0) {$V_2$};
\node[font=\small,align=center] at (1.7,-1.3)
  {依原卷关系重绘的示意图；非比例图。};
\end{tikzpicture}
\end{center}
\nopagebreak[4]
\Needspace{8\baselineskip}

\Needspace{4\baselineskip}
\textbf{解析}

\gksubpart{（1）}

双曲线中 $a=4,b=3,c=5,e=\dfrac54$，准线为 $x=\pm\dfrac ae=\pm\dfrac{16}{5}$，两条都须写出。

\gksubpart{（2）}

展开各项，$x^2$ 项系数为 $-1-3-6=-10$。

\gksubpart{（3）}

令公差为 $d\ne0$。当 $n\to\infty$ 时，$a_n\sim nd$、$S_n\sim\dfrac12dn^2$，故 $\dfrac{na_n}{S_n}\to2$。

\gksubpart{（4）}

$\sin x+\cos x=\sqrt2\sin(x+\frac\pi4)$，最大值为 $\sqrt2$。

\gksubpart{（5）}

以棱柱底面为截面。分割平面一侧的三角形截面积，相对底面积随棱柱高度参数 $t\in[0,1]$ 按 $\dfrac{(1+t)^2}{4}$ 变化，故该侧体积分数为 $\int_0^1\dfrac{(1+t)^2}{4}\,dt=\dfrac7{12}$；另一侧为 $\dfrac5{12}$，比为 $7:5$。该结果也与来源答案一致。

\gkanswer{答案（1）：$x=\pm\dfrac{16}{5}$}

\gkanswer{答案（2）：$-10$}

\gkanswer{答案（3）：$2$}

\gkanswer{答案（4）：$\sqrt2$}

\gkanswer{答案（5）：$7:5$}

\end{answer}
\begin{answer}
\gksolheadwithlabel{27}{二}{}

\textbf{题目}：

（本题满分 13 分）

如图，在三棱锥 $S-ABC$ 中，$SA\perp$ 底面 $ABC$，$AB\perp BC$，$DE$ 垂直平分 $SC$，且分别交 $AC,SC$ 于 $D,E$，又 $SA=AB=a$，$BC=\sqrt2a$。

（1） 求证：$SC\perp$ 面 $BDE$；

（2） 求面 $BDE$ 与面 $BDC$ 所成二面角的大小。
\nopagebreak[4]
\begin{center}\begin{tikzpicture}[x=0.95cm,y=0.78cm,line cap=round,line join=round]
\coordinate (A) at (-1.6,0);
\coordinate (C) at (2.4,0);
\coordinate (B) at (0.4,-1.65);
\coordinate (S) at (-1.6,3.1);
\coordinate (D) at (-0.8,0);
\coordinate (E) at (0.4,1.55);

% 依原卷中点、垂直平分线及棱锥的关系重绘；不表达额外长度比例。
\draw[thick] (A)--(S)--(C)--(B)--cycle;
\draw[thick] (S)--(B);
\draw[dashed] (A)--(C) (D)--(B);
\draw[dashed,thick] (D)--(E);
\draw[thick] (S)--(E);
\draw[thick] (E)--(C);

% SC 的等长刻痕。D、E 的投影位置保留原卷的垂直平分线关系。
\draw (-0.71,2.185) -- (-0.49,2.465);
\draw (1.29,0.635) -- (1.51,0.915);
% DE 与 SC 的直角标记，按两边方向构造小方角。
\draw (0.266,1.376) -- (0.44,1.242) -- (0.574,1.416);

\node[above left] at (S) {$S$};
\node[left] at (A) {$A$};
\node[below] at (B) {$B$};
\node[right] at (C) {$C$};
\node[below] at (D) {$D$};
\node[above right] at (E) {$E$};
\node[font=\small,align=center] at (0.4,-2.35)
  {依原卷关系重绘的示意图；非比例图。};
\end{tikzpicture}
\end{center}
\nopagebreak[4]
\Needspace{4\baselineskip}
\textbf{解析}

\gksubpart{（1）}

因 $SA\perp$ 底面 $ABC$，故 $SA\perp AB$。由 $SA=AB=a$，直角三角形 $SAB$ 中 $SB=\sqrt2a=BC$。又 $E$ 是 $SC$ 中点，等腰三角形 $SBC$ 的中线 $BE$ 垂直于 $SC$。已知 $DE\perp SC$，且 $BE,DE$ 相交于 $E$ 并在平面 $BDE$ 内，因此 $SC\perp$ 面 $BDE$。

\gksubpart{（2）}

由（1）得 $SC\perp BD$；又 $SA\perp$ 底面 $ABC$，故 $SA\perp BD$。因 $SC$ 与 $SA$ 相交于 $S$ 且同在平面 $SAC$ 内，得 $BD\perp$ 平面 $SAC$。两平面与 $SAC$ 的交线分别为 $DE$、$DC$，故平面角为 $\angle EDC$。在直角三角形 $ABC$ 中，$AC=\sqrt3a$；且 $SA\perp AC$、$SA=a$，所以 $\tan\angle ACS=1/\sqrt3$，从而 $\angle ACS=30^\circ$。结合 $DE\perp SC$，得 $\angle EDC=60^\circ$。

\gkanswer{答案（1）：见解析版}

\gkanswer{答案（2）：$60^\circ$}

\end{answer}
\begin{answer}
\gksolheadwithlabel{28}{三}{}

\textbf{题目}：

（本题满分 14 分）

已知椭圆的焦点为 $F_1(0,-1)$ 和 $F_2(0,1)$，直线 $y=4$ 是椭圆的一条准线。

（1） 求椭圆的方程；

（2） 又设点 $P$ 在这个椭圆上，且 $|PF_1|-|PF_2|=1$，求 $\operatorname{tg}\angle F_1PF_2$ 的值。

\Needspace{4\baselineskip}
\textbf{解析}

\gksubpart{（1）}

设椭圆长半轴为 $a$、短半轴为 $b$，半焦距 $c=1$。准线到中心的距离为 $a^2/c=4$，故 $a^2=4$、$b^2=a^2-c^2=3$，方程为 $\dfrac{y^2}{4}+\dfrac{x^2}{3}=1$。

\gksubpart{（2）}

椭圆定义给出 $|PF_1|+|PF_2|=2a=4$，与差为 $1$ 联立得 $|PF_1|=\dfrac52$、$|PF_2|=\dfrac32$。又 $|F_1F_2|=2$，余弦定理得 $\cos\angle F_1PF_2=\dfrac35$，因而 $\operatorname{tg}\angle F_1PF_2=\dfrac43$。

\gkanswer{答案（1）：$\dfrac{y^2}{4}+\dfrac{x^2}{3}=1$}

\gkanswer{答案（2）：$\operatorname{tg}\angle F_1PF_2=\dfrac43$}

\end{answer}
\begin{answer}
\gksolheadwithlabel{29}{四}{}

\textbf{题目}：

（本题满分 14 分）

\Needspace{5\baselineskip}
已知数列 $\{a_n\}$ 的前 $n$ 项和的公式是
\nopagebreak[4]
\[
S_n=\frac\pi{12}(2n^2+n)
\]
。

（1） 求证 $\{a_n\}$ 是等差数列，并求出它的首项和公差；

\Needspace{5\baselineskip}
（2） 记 $b_n=\sin a_n\cdot\sin a_{n+1}\cdot\sin a_{n+2}$，求证：对任何自然数 $n$，都有
\nopagebreak[4]
\[
b_n=\frac{\sqrt2}{8}(-1)^{n-1}
\]
。

\Needspace{4\baselineskip}
\textbf{解析}

\gksubpart{（1）}

$a_1=S_1=\dfrac\pi4$。当 $n\ge2$ 时，$a_n=S_n-S_{n-1}=\dfrac\pi{12}(4n-1)$；该式对 $n=1$ 也成立。因此 $a_{n+1}-a_n=\dfrac\pi3$，数列为等差数列，首项 $\dfrac\pi4$，公差 $\dfrac\pi3$。

\gksubpart{（2）}

由通项 $a_n=\dfrac\pi{12}(4n-1)$ 可知 $a_n/\pi$ 不为整数，故 $\sin a_n\ne0$。由公差得 $a_{n+3}=a_n+\pi$，于是 $\dfrac{b_{n+1}}{b_n}=\dfrac{\sin a_{n+3}}{\sin a_n}=-1$。又 $b_1=\sin\dfrac\pi4\sin\dfrac{7\pi}{12}\sin\dfrac{11\pi}{12}=\dfrac{\sqrt2}{8}$，所以 $b_n=\dfrac{\sqrt2}{8}(-1)^{n-1}$。

\gkanswer{答案（1）：$a_1=\dfrac\pi4$，$d=\dfrac\pi3$}

\gkanswer{答案（2）：$b_n=\dfrac{\sqrt2}{8}(-1)^{n-1}$}

\end{answer}
\begin{answer}
\gksolheadwithlabel{30}{五}{}

\textbf{题目}：

（本题满分 14 分）

\Needspace{5\baselineskip}
设
\nopagebreak[4]
\[
f(x)=\lg\frac{1+2^x+4^xa}{3},\qquad a\in\mathbb R
\]
。

（1） 如果 $0<a\le1$，求证：当 $x\ne0$ 时，有 $2f(x)<f(2x)$；

（2） 如果当 $x\in(-\infty,1]$ 时 $f(x)$ 有意义，求 $a$ 的取值范围。

\begin{gktip}
完整试题共 30 个内部题目块。第 12 题图位缺失状态保持待核；本文件仅为来源转录，来源答案与独立解析另行登记。
\end{gktip}

\Needspace{4\baselineskip}
\textbf{解析}

\gksubpart{（1）}

\Needspace{5\baselineskip}
令 $p=1,q=2^x,r=4^xa$。因 $x\ne0$，$p\ne q$，有 $2pq<p^2+q^2$；另有 $2qr\le q^2+r^2$、$2rp\le r^2+p^2$。将三式相加，得 $(p+q+r)^2<3(p^2+q^2+r^2)$。又 $0<a\le1$ 给出 $a^2\le a$，故
\nopagebreak[4]
\[
\frac{(1+2^x+4^xa)^2}{9}<\frac{1+2^{2x}+4^{2x}a}{3}
\]
。
两边均为正，取常用对数即得 $2f(x)<f(2x)$。

\gksubpart{（2）}

\Needspace{5\baselineskip}
定义域要求对所有 $x\le1$ 有 $1+2^x+4^xa>0$。除以正数 $4^x$，等价于
\nopagebreak[4]
\[
a>-\left(\left(\frac14\right)^x+\left(\frac12\right)^x\right)\qquad(x\le1)
\]
。
右侧随 $x$ 增大而增大，在 $x=1$ 取得最大值 $-\dfrac34$，所以充要条件为 $a>-\dfrac34$。

\gkanswer{答案（1）：见解析版}

\gkanswer{答案（2）：$a>-\dfrac34$}

\end{answer}